% generated by agent/render_tools.py from the server's list_tools
\begin{longtable}{@{}l p{0.3\linewidth} p{0.45\linewidth}@{}}
\toprule tool & arguments & what it does\\ \midrule \endhead
\texttt{camera\_info} & -- & Describe the camera: vendor, model, sensor size, bit depth (max DN), exposure and gain ranges, the lab limits, and the current settings (exposure, gain, ROI).\\
\texttt{set\_exposure} & exposure\_us: number & Set the exposure time in microseconds.\\
\texttt{set\_gain} & gain\_db: number & Set the analog gain in dB.\\
\texttt{set\_roi} & x: integer, y: integer, width: integer, height: integer & Set the region of interest: offset (x, y) and size in sensor pixels.\\
\texttt{reset\_roi} & -- & Use the full sensor again.\\
\texttt{set\_illumination} & level: number & Light source level 0..1 on the simulated camera only (0 = dark, for dark frames).\\
\texttt{capture} & n\_frames: integer (opt.), label: string (opt.), regions: array (opt.) & Acquire n\_frames frames with the current settings and save them.\\
\texttt{capture\_series} & n\_frames: integer, interval\_s: number, label: string (opt.), regions: array (opt.) & Stability measurement: n\_frames frames, one every interval\_s seconds, with the current settings.\\
\texttt{analyze\_frame} & frame\_file: string, black\_level\_dn: number (opt.), region: object (opt.) & Speckle and beam measures of a saved frame (first frame of the file): speckle contrast and grain size (FWHM of the intensity autocorrelation) inside `region` (or the whole frame), and the beam centroid and D4sigma widths.\\
\bottomrule
\end{longtable}
